% Cover letter for the Journal of Informetrics submission of Paper E. Numbers match content/abstract.tex and FINDINGS.md.
\documentclass[10pt]{letter}
\usepackage[margin=1.6cm,top=1.2cm,bottom=1.4cm]{geometry}
\usepackage{xurl}
\sloppy
\usepackage[hidelinks]{hyperref}
\signature{Hossein Mohasel Arjomandi\\Microsoft, Redmond, WA, USA\\\href{mailto:hosseinm@microsoft.com}{hosseinm@microsoft.com}}
\address{}
\date{\today}
\begin{document}
\begin{letter}{Editorial Office\\\emph{Journal of Informetrics}}
\opening{Dear Editor,}

I am pleased to submit the manuscript \emph{``Do retracted papers look disruptive? The CD index before retraction on
the open biomedical citation graph''} for consideration as a research article in the \emph{Journal of Informetrics}.

\textbf{Context.} Proposals to use the CD index in research evaluation need evidence on how it treats unreliable
work. This study measures papers that were later retracted, with scores computed before their retraction.

\textbf{What the manuscript does.} Using only open data (OpenCitations, January 2024, linked to PubMed and to
Retraction Watch), I computed the five-year CD index for 4,587 biomedical papers retracted only after their scoring
window had closed, comparing 4,586 of them with 31,662 scored, never-retracted controls in 2,239 journal--year strata.
The design was specified before scoring; frozen-score hashes were checked before labels were joined. The findings are:
\begin{itemize}
\item The mean score difference is small: a post hoc 90\% interval rules out absolute differences above 0.00093,
  about 28\% of the review comparison's effect. In the pre-specified five-year analysis, there is no clear over-representation among the top 5\%
  (odds ratio 0.99, 95\% CI 0.83--1.20), but positive scores are more common (1.23, 1.13--1.34).
\item Post hoc analyses apply a known denominator problem to retracted papers: their citing papers are less likely
  to also cite their references (65.5\% versus 68.8\%). The variant without $n_k$ shows an unfiltered mean gap
  of $+0.067$ (95\% whole-journal bootstrap CI 0.034--0.095).
  Removing notice, comment/letter and retracted citing records is a robustness check and raises the gap to $+0.082$.
  The contrast persists with reference and citation adjustment and after excluding the ten highest-retraction journals.
  The post hoc three-year window leans toward a modest top-tier excess, though this weakens under author-name exclusion.
\item Post-retraction scores are shifted by notice records and a widening citation-composition gap, partly continuing
  an earlier trend; scoring windows should end before the notice.
\end{itemize}

\textbf{Why the journal.} The study contributes to the journal's debate on disruption-index validity and data
dependence. Added analyses are labelled post hoc and were planned before execution. No paper-level scores of
non-retracted papers are released.

\textbf{Openness.} Open inputs, code, specifications, checksummed run records and a claim-to-run map are documented at
\mbox{\href{https://github.com/HosseinMA96/retracted-disruption-open-data}{github.com/HosseinMA96/retracted-disruption-open-data}}.

\textbf{Related submissions.} Two other manuscripts by the same author are under consideration at
\emph{Scientometrics}: \emph{``Can open citation data detect citation stacking? An audit against Clarivate's journal
suppressions''} and \emph{``Do links to retracted papers predict future retractions? An out-of-sample test on the
open biomedical citation graph''}; the latter is also available as the preprint
\href{https://arxiv.org/abs/2611.10474}{arXiv:2611.10474}. This manuscript reuses their open citation snapshot and the
second study's linked Retraction Watch dataset
(\href{https://doi.org/10.5281/zenodo.23874477}{doi:10.5281/zenodo.23874477}).
The companion preprint tests whether links to already-retracted papers predict future retractions.
This manuscript instead examines disruption-index scores before a paper's first retraction notice and reports none
of the companion manuscripts' results.

The manuscript is original, has not been published, and is not under consideration elsewhere. I am the sole author,
there is no funding to declare and I have no competing interests. The use of generative AI is declared in line with
Elsevier's policy and described in Methods. The views expressed are my
own and do not necessarily reflect those of Microsoft.

Thank you for considering this manuscript.

\closing{Sincerely,}
\end{letter}
\end{document}
